\subsection{CROSSES}

We have also defined in algebra the cross of a pair of lines in a two-dimensional vector space over a maximal ordered field (*). This definition applies in particular to the real plane \(\mathbf{R}^2\) . The set \(\mathfrak{A}_0\) of all crosses has the structure of an abelian group (written additively) defined by

\[
(\widehat {\mathbf {D _ {1} , D _ {3}}}) = (\widehat {\mathbf {D _ {1} , D _ {2}}}) + (\widehat {\mathbf {D _ {2} , D _ {3}}})
\]

so that, in particular, \((\widehat{\mathbf{D}_1,\mathbf{D}_1}) = 0\) and \((\widehat{\mathbf{D}_2,\mathbf{D}_1}) = -(\widehat{\mathbf{D}_1,\mathbf{D}_2}).\)

The right cross \(\delta_0\) is the solution \(\neq 0\) of the equation \(2\theta = 0\) in \(\mathfrak{A}_0\) ; it is the cross which the imaginary axis makes with the real axis.

We define a canonical homomorphism \(\varphi\) of the group \(\mathfrak{A}\) of angles onto the group \(\mathfrak{A}_0\) of crosses by making correspond to the angle which a ray \(\Delta\) makes with the positive real semi-axis, the cross which the line D containing \(\Delta\) makes with the real axis. A cross \(\theta_0\) is the image under \(\varphi\) of two angles \(\theta\) and \(\theta + \varpi\) ; thus \(\mathfrak{A}_0\) is isomorphic to the quotient of \(\mathfrak{A}\) by the subgroup \(\{o, \varpi\}\) . If we transport to \(\mathfrak{A}_0\) the topology of the quotient group \(\mathfrak{A}/\{o, \varpi\}\) by means of the bijective homomorphism associated with \(\varphi\) , then \(\mathfrak{A}_0\) becomes a compact topological group and \(\varphi\) a strict morphism of \(\mathfrak{A}\) onto \(\mathfrak{A}_0\) .

If we compose the homomorphism \(\varphi\) of \(\mathfrak{A}\) onto \(\mathfrak{A}_0\) with the homomorphism \(x \to \Im(x/a)\) of \(\mathbf{R}\) onto \(\mathfrak{A}\) , we have a homomorphism \(x \to \Im_0(x/a)\) of \(\mathbf{R}\) onto \(\mathfrak{A}_0\) ; every cross \(\theta_0 \in \mathfrak{A}_0\) corresponds, under this homomorphism, to a class of real numbers mod \(\frac{1}{2}a\) , which is called the measure of the cross \(\theta_0\) (relative to the base \(a\) ); by abuse of language, every number in this class is called a measure of \(\theta_0\) , and that one which belongs to the interval \([0, \frac{1}{2}a]\) is called the principal measure of \(\theta_0\) ; the cross \(\Im_0(x/a)\) is the cross of measure \(x\) . Every measure of \(\theta_0\) is also a measure of one of the two angles \(\theta, \theta + \varpi\) whose image under the homomorphism \(\varphi\) is \(\theta_0\) .

Here again, once the base a has been chosen, when we speak of a cross we generally mean, by abuse of language, a measure of this cross relative to the base a.

Remark. We can define a homomorphism of \(\mathbf{C}^*\) onto \(\mathfrak{A}_0\) by mapping each complex number \(z \neq 0\) to the cross which the line through \(0\) and \(z\) makes with the real axis. Clearly this homomorphism is the composition of \(\varphi\) and the homomorphism \(z \to \operatorname{Am}(z)\) of \(\mathbf{C}^*\) onto \(\mathfrak{A}\) ; it is therefore a strict morphism of the topological group \(\mathbf{C}^*\) onto the topological group \(\mathfrak{A}_0\) , and the associated bijective homomorphism is an isomorphism of the quotient group \(\mathbf{C}^*/\mathbf{R}^*\) onto \(\mathfrak{A}_0\) .

We know that if D denotes a line making a cross \(\theta_0\) with the real axis and if \((a, b)\) is a pair of direction ratios of D, then the tangent of the cross \(\theta_0\) (denoted by tan \(\theta_0\) ) is the element \(b/a\) of \(\tilde{\mathbf{R}} (= \infty\) if \(a = 0\) ), which is also called the slope of the line D. If \(\theta\) and \(\theta + \varpi\) are the two angles whose image under \(\varphi\) is \(\theta_0\) , then we have \(\tan \theta_0 = \tan \theta = \tan (\theta + \varpi)\) . The mapping \(\theta_0 \to \tan \theta_0\) is a homeomorphism of \(\mathfrak{A}_0\) onto \(\tilde{\mathbf{R}}\) , for the topological space \(\mathbf{C}^*/\mathbf{R}^*\) is just the real projective line \(\mathbf{P}_1\) , and from Chapter VI, § 3, no. 3, the mapping of a line (considered as a point of \(\mathbf{P}_1\) ) to its slope is a homeomorphism of \(\mathbf{P}_1\) onto \(\tilde{\mathbf{R}}\) . If now we transport to \(\tilde{\mathbf{R}}\) the group structure of \(\mathfrak{A}_0\) by means of the mapping \(\theta_0 \to \tan \theta_0\) , we have defined on \(\tilde{\mathbf{R}}\) the structure of an abelian topological group, in which the product of two elements \(t_1\) , \(t_2\) is \(\frac{t_1 + t_2}{1 - t_1 t_2}\) whenever \(t_1\) , \(t_2\) belong to R and \(t_1 t_2 \neq 1\) ; for pairs \((t_1, t_2)\) which do not satisfy these conditions, the product of \(t_{1}\) and \(t_{2}\) is obtained by extending the function \(\frac{x+y}{1-xy}\) by continuity to \(\tilde{R}\times\tilde{R}\) , and is still denoted by \(\frac{t_{1}+t_{2}}{1-t_{1}t_{2}}\) .

